% Section: Theorem K -- sums of k infinite sets in countable abelian groups.
% Label prefix: K:
\section{Sums of \texorpdfstring{$k$}{k} infinite sets in countable abelian groups}
\label{K:sec:main}

This section answers Question~5.14 of~\cite{kra2023problems}.
We show (\cref{K:thm:main}) that in every countably infinite abelian group, every set of
positive upper Banach density contains a sumset $B_1+\dots+B_k$ of $k$ infinite sets, for
every $k$.
For finite groups the question has a trivially negative answer (\cref{K:rem:finite}), so
``countably infinite'' is the right setting.
The proof does not use higher-order Host--Kra theory, nilsystems or nilspaces, or any
equidistribution theorem for nilsystems; besides the Kronecker factor, it uses only standard
facts of ergodic and spectral theory.
It extends to $k$ summands, and to arbitrary countable abelian groups, the short proof of
the case $k=2$ for subsets of $\N$ given by Kra, Moreira, Richter and
Robertson~\cite{kra2026short}.
The cube is built one level at a time.
At each level we re-base on the Kronecker factor of the current measure, and from one level
to the next we keep only a condition on topological supports.

\subsection{The question and the result}\label{K:ssec:result}

Let $\Gamma$ be a countable abelian group, written additively.
A \emph{\Fol{} sequence} in $\Gamma$ is a sequence $\Phi=(\Phi_N)_{N\in\N}$ of nonempty finite
subsets of $\Gamma$ such that $|(\Phi_N+\gamma)\triangle\Phi_N|/|\Phi_N|\to0$ for every
$\gamma\in\Gamma$.
The \emph{upper Banach density} of $A\subset\Gamma$ is
\[
\ubd(A)=\sup_{\Phi}\ \limsup_{N\to\infty}\frac{|A\cap\Phi_N|}{|\Phi_N|},
\]
where the supremum is over all \Fol{} sequences $\Phi$ in $\Gamma$.

\begin{question}{5.14}\label{K:q:514}
Is it true in every countable abelian group $(G,+)$ that every set $A$ of positive upper
Banach density contains $B_1+\dots+B_k$ for infinite sets $B_1,\dots,B_k\subset G$?
\end{question}

\begin{theorem}\label{K:thm:main}
Let $\Gamma$ be a countably infinite abelian group and let $k\ge1$.
If $A\subset\Gamma$ satisfies $\ubd(A)>0$, then there are infinite sets
$B_1,\dots,B_k\subset\Gamma$ with
\[
B_1+\dots+B_k=\{b_1+\dots+b_k:\ b_i\in B_i\}\subset A .
\]
\end{theorem}

\begin{remark}[Finite groups]\label{K:rem:finite}
Read literally, \cref{K:q:514} includes finite groups, and for them the answer is no.
A finite group has no infinite subsets, yet every nonempty $A\subset G$ has
$\ubd(A)\ge|A|/|G|>0$.
(For finite $G$, the Følner condition forces $\Phi_N=G$ for all large $N$.)
So the meaningful version of the question concerns countably \emph{infinite} groups, and
\cref{K:thm:main} answers that version positively for every $k$.
\end{remark}

\begin{remark}[Previously known cases]\label{K:rem:prior}
\leavevmode
\begin{enumerate}[label=(\alph*)]
\item For $k=2$ the conclusion holds in every countably infinite amenable group, with products in
place of sums, for sets of positive density along a two-sided \Fol{} sequence. This is due
to Moreira, Richter and Robertson~\cite[Theorem~1.3]{moreira2019proof}. In an abelian group
every \Fol{} sequence is two-sided.
\item For $\Gamma=\Z$, or rather for $A\subset\N$, the case of general $k$ is
\cite[Theorem~1.1]{kra2022infinite} (Theorem~2.8 of~\cite{kra2023problems}). Here and below,
theorem numbers in~\cite{kra2022infinite} refer to arXiv:2206.01786v2.
Its proof relies on pronilfactors: Kra, Moreira, Richter and Robertson construct Erdős
cubes in an extension of the system with continuous factor maps to its pronilfactors
\cite[Section~5]{kra2022infinite}.
\item For $k=3$ and $[\Gamma:6\Gamma]<\infty$, the conclusion follows from a theorem of Ackelsberg
and Jamneshan~\cite[Theorem~1.9]{ackelsberg2025equidistribution}. They find $t\in\Gamma$ and an
infinite $B\subset\Gamma$ such that $A-t$ contains every sum of three distinct elements of
$B$. Splitting $B$ into three disjoint infinite sets $B'_1,B'_2,B'_3$ then gives
$(B'_1+t)+B'_2+B'_3\subset A$.
\end{enumerate}
After stating \cref{K:q:514}, the survey~\cite{kra2023problems} remarks that the techniques
of~\cite{kra2022infinite}, together with work of Griesmer~\cite{griesmer2009ergodic}, may
suffice for $\Gamma=\Z^d$. In other groups, it says, ``one cannot
proceed in the same way, because the analogous ergodic structure theory is not currently
available.''
The proof below uses no such structure theory.
\end{remark}

\begin{remark}[The integers]\label{K:rem:N}
Our proof specialises to a proof of \cite[Theorem~1.1]{kra2022infinite} that uses no
pronilfactors.
To see this, let $A\subset\N$ have positive upper Banach density.
Intervals form a \Fol{} sequence in $\Z$, so $\ubd(A)>0$ also when $A$ is viewed as a subset
of $\Z$.
\cref{K:thm:main} then gives infinite $B_i\subset\Z$ with $B_1+\dots+B_k\subset A\subset\N$.
Fix $b_j\in B_j$ for $j\neq i$. Then every element of $B_i$ is at least
$1-\sum_{j\ne i}b_j$, so each $B_i$ is bounded below.
Hence $B_i\cap\N$ is still infinite, and $(B_1\cap\N)+\dots+(B_k\cap\N)\subset A$.
\end{remark}

\subsection{Strategy of the proof}\label{K:ssec:strategy}

\subsubsection*{Erdős cubes}
Following \cite{kra2022infinite}, we look for a dynamical configuration called an Erdős cube.
The correspondence principle (\cref{K:thm:corr}) gives the following data:
\begin{itemize}
\item an ergodic system $(X,\mu,T)$;
\item an open set $E\subset X$ with $\mu(E)>0$;
\item a point $a\in X$, generic for $\mu$, such that $A=\{\gamma:T^\gamma a\in E\}$.
\end{itemize}
An \emph{Erdős cube} is a tuple $y=(y_\eps)_{\eps\in\{0,1\}^k}$ of points of $X$ with the
following property. For every direction $i$, the lower $i$-face
$(y_\eps)_{\eps_i=0}$ is moved arbitrarily close to the upper $i$-face
$(y_{\eps+e_i})_{\eps_i=0}$ by infinitely many group elements $\beta$, acting diagonally;
here $\eps+e_i$ is $\eps$ with its $i$-th digit changed to $1$.
If $y_{0\cdots0}=a$ and $y_{1\cdots1}\in E$, a greedy argument produces
$B_1+\dots+B_k\subset A$ (\cref{K:lem:cube-sumset}).
So the task is to construct such a cube.

\subsubsection*{The case $k=2$}
The argument of~\cite{kra2026short} can be read as constructing a $2$-dimensional cube
whose bottom face is $(a,y)$ and whose top face $(x,x')$ is random; we use our own
notation. A point $y\in\supp\mu$ is chosen so that $(a,y)$
is generic for some joining $\lambda$ of $\mu$ with itself.
The key claim (Claim~1 in the proof of \cite[Theorem~3.1]{kra2026short}) produces
infinitely many $\beta$ with two properties:
\begin{itemize}
\item $T^\beta a$ is close to $y$;
\item the top face $(x,x')$, which is distributed according to $\lambda$, still has positive
measure after the direction-$1$ constraint imposed by $\beta$.
\end{itemize}
The proof decomposes $1_E=f_c+f_{wm}$ into its Kronecker and weak mixing parts.
\begin{itemize}
\item The weak mixing part averages out along every \Fol{} sequence.
\item The Kronecker part is handled by approximating $f_c$ by \emph{continuous}
eigenfunctions. This uses an enhanced correspondence principle in which every measurable
eigenfunction is topological \cite[Theorem~2.5(4)]{kra2026short}. Continuity gives the
equicontinuity estimate~(4.2) of~\cite{kra2026short}, which turns ``$T^\beta a$ is close to
$y$'' into ``$f_c\circ T^\beta$ is close to $f_c$ along the orbit''.
\item The positive frequency of the $\beta$ with $T^\beta a$ near $y$ comes from
$y\in\supp\mu$.
\end{itemize}

\subsubsection*{Where the difficulty lies for $k\ge3$}
Suppose a $(k-1)$-dimensional bottom face $w=(w_\eps)$ with $w_{0\cdots0}=a$ has been
built, generic for an ergodic measure $\lambda$ on $X^{\{0,1\}^{k-1}}$. We try to add a
random top face distributed according to $\lambda$.
\begin{itemize}
\item The direction $k$ (bottom face to top face) behaves exactly as in the case $k=2$.
\item In each of the other directions $i$, the relevant system is the face measure $\nu_i$,
the distribution under $\lambda$ of the upper $i$-face. Its Kronecker factor carries
information about $\mu$ beyond the Kronecker factor of $\mu$. For instance, for the skew
product $(s,u)\mapsto(s+\theta,u+s)$ on $\T^2$, the measures on pairs that arise already
have eigenfunctions of the form $e(u_1-u_0)$, with eigenvalues that depend on the
parameter of the pair measure. For a general system there is no reason for such
eigenfunctions to be continuous in the topological model in which $a$ lives.
\end{itemize}
So the equicontinuity step of~\cite{kra2026short} is not available on faces.
The proof of~\cite{kra2022infinite} for $\Z$ proceeds differently. It disintegrates the
Host--Kra cubic measure $\mu^{[k]}$ over the first coordinate and needs the conditional
measure at the prescribed point $a$. To define that conditional measure at every point,
it passes to an extension in which the factor maps onto the pronilfactors are continuous,
and it proves that the ergodic decomposition of $\mu^{[k]}$ is then continuous.

\subsubsection*{The three ideas}
\begin{enumerate}[label=(\roman*)]
\item \textbf{Phase returns replace continuity} (\cref{K:prop:criterion}).
Let $\iota_i$ be the Kronecker map of $\nu_i$, with rotation $\alpha_i$.
\begin{itemize}
\item By ergodicity, $\iota_i(\text{upper face})-\iota_i(\text{lower face})$ is a
$\lambda$-a.e.\ constant $r_i$.
\item By continuity of translation in $L^1$ of the Kronecker group, the Kronecker part
is controlled for all $\beta$ whose ``phase'' $\alpha_i(\beta)$ is close to $r_i$.
\item So it suffices to have many $\beta$ for which $T^\beta(\text{lower face of }w)$ is close
to the upper face of $w$ \emph{and} $\alpha_i(\beta)$ is close to $r_i$. Here ``many'' means
a set of positive upper Banach density.
\end{itemize}
\item \textbf{Phase returns come from supports} (\cref{K:lem:phase}).
Such phase returns exist as soon as the upper face of $w$ lies in the topological support
of $\nu_i$. The proof re-bases $w$: we change the \Fol{} sequence so that $w$, decorated
with a point of the Kronecker group, becomes generic for the graph measure. A stabiliser
argument for the support of the graph measure then forces the decorated upper face into
that support (\cref{K:lem:stab}).
\item \textbf{The support condition survives one more level} (\cref{K:prop:step}).
We build the cube level by level. At level $m$ we have a point $w\in X^{\{0,1\}^m}$ that
is generic for an ergodic, face-symmetric measure $\lambda_m$ and whose upper faces lie in
the supports of the face measures. To pass to level $m+1$:
\begin{itemize}
\item re-base $w$ on the Kronecker factor of $\lambda_m$;
\item choose the new top face at random, so that the new cube is generic for an ergodic
component of the relative product of $\lambda_m$ with itself over its Kronecker factor;
\item check that the support condition persists. This is the content of an elementary
convolution lemma (\cref{K:lem:PS}). The lemma holds at \emph{every} support point, not
only at typical ones.
\end{itemize}
\end{enumerate}
From one level to the next, the induction carries only genericity, ergodicity, face
symmetry, the vertex marginals and the support condition; the Kronecker data are rebuilt at
each level.
The only prescribed point is $a$. Every vertex added at an induction step is a typical
point for the current measure, chosen after all the data it must avoid have been fixed, and
the final top face is constructed by nested boxes. Apart from genericity, no statement
about typical points is ever applied at a prescribed point, and this is why no
equidistribution at a prescribed point is needed.

\subsection{Notation and standard facts}\label{K:ssec:prelim}

Throughout, $\Gamma$ is a countably infinite abelian group.
For a finite set $F$ and a function $u$ on $F$ we write
$\mathop{\mathbb E}_{\gamma\in F}u(\gamma)=|F|^{-1}\sum_{\gamma\in F}u(\gamma)$.

\begin{lemma}\label{K:lem:folner-grow}
If $\Phi$ is a \Fol{} sequence in $\Gamma$, then $|\Phi_N|\to\infty$.
\end{lemma}

\begin{proof}
Let $F\subset\Gamma$ be finite with $|F|=n$. The number of pairs $(x,g)\in\Phi_N\times F$
with $x+g\in\Phi_N$ is at most $|\Phi_N|^2$. It also equals
$\sum_{g\in F}|\Phi_N\cap(\Phi_N-g)|$, which is $n|\Phi_N|(1-o(1))$ by the \Fol{}
property. Hence $|\Phi_N|\ge n(1-o(1))$. Since $\Gamma$ is infinite, $n$ is arbitrary.
\end{proof}

\subsubsection*{Topological systems and generic points}
A \emph{topological $\Gamma$-system} $(Y,T)$ consists of a compact metric space $Y$ and an
action of $\Gamma$ by homeomorphisms $T^\gamma$, with $T^{\gamma+\delta}=T^\gamma T^\delta$.
All measures are Borel probability measures.
A measure $\nu$ is \emph{invariant} if $T^\gamma_*\nu=\nu$ for all $\gamma$, and
\emph{ergodic} if moreover every invariant Borel set has measure $0$ or $1$.
Since $\Gamma$ is countable, this is equivalent to asking that every Borel set $S$ with
$\nu(S\triangle T^{-\gamma}S)=0$ for all $\gamma$ has measure $0$ or $1$ (replace $S$ by
$\bigcap_{\gamma}T^{-\gamma}S$); in particular a Borel map into a compact metric space
which is $\nu$-a.e.\ invariant under each $T^\gamma$ is $\nu$-a.e.\ constant.
A point $y\in Y$ is \emph{generic for $\nu$ along $\Phi$}, written $y\in\mathrm{gen}(\nu,\Phi)$,
if $\mathop{\mathbb E}_{\gamma\in\Phi_N}f(T^\gamma y)\to\int f\,d\nu$ for every $f\in C(Y)$.
The (topological) support $\supp\nu$ is the smallest closed set of full measure. It is
invariant whenever $\nu$ is.
We use the following facts freely.
\begin{enumerate}[label=(G\arabic*)]
\item\label{K:G1} If $y\in\mathrm{gen}(\nu,\Phi)$ and $O\subset Y$ is open, then
$\liminf_N|\{\gamma\in\Phi_N:T^\gamma y\in O\}|/|\Phi_N|\ge\nu(O)$.
This follows by approximating $1_O$ from below by continuous functions.
In particular, if $O$ meets $\supp\nu$, then the set of return times
$\{\gamma:T^\gamma y\in O\}$ has positive upper Banach density. By \cref{K:lem:folner-grow}
it is then infinite.
\item\label{K:G2} If $y\in\mathrm{gen}(\nu,\Phi)$, then $\supp\nu$ is contained in the orbit closure
$\overline{\{T^\gamma y:\gamma\in\Gamma\}}$. This is immediate from \ref{K:G1}.
\item\label{K:G3} If $P\colon Y\to Y'$ is continuous and equivariant and
$y\in\mathrm{gen}(\nu,\Phi)$, then $P(y)\in\mathrm{gen}(P_*\nu,\Phi)$.
\item\label{K:G4} If $\nu$ is ergodic, then for every \Fol{} sequence $\Phi$ there is a
subsequence $\Psi$ such that $\nu$-almost every point is generic for $\nu$ along $\Psi$
\cite[Lemma~2.7]{charamaras2024finding}. This rests on Lindenstrauss' pointwise ergodic
theorem.
\end{enumerate}

\subsubsection*{Cubes}
For $m\ge0$ let $[m]=\{0,1\}^m$, with elements written $\eps=\eps_1\cdots\eps_m$.
Let $\vec0=0\cdots0$ and $\vec1=1\cdots1$.
$\Gamma$ acts on $X^{[m]}$ diagonally, $(T^\gamma x)_\eps=T^\gamma x_\eps$; we use the same
letter $T$ for this action.
We write $e_i\in[m]$ for the vertex with a single digit $1$, in position $i$, and
$\eps+e_i$ for the vertex obtained from $\eps$ (with $\eps_i=0$) by changing its $i$-th
digit to $1$.
For $1\le i\le m$, $\eta\in[m-1]$ and $j\in\{0,1\}$, let $\eta^{(i,j)}\in[m]$ be the vertex
obtained by inserting the digit $j$ at position $i$.
The \emph{lower} and \emph{upper $i$-faces} are the continuous equivariant maps
$L_i,U_i\colon X^{[m]}\to X^{[m-1]}$ given by
\[
(L_ix)_\eta=x_{\eta^{(i,0)}},\qquad (U_ix)_\eta=x_{\eta^{(i,1)}} .
\]
We identify $X^{[m+1]}$ with $X^{[m]}\times X^{[m]}$ by $y\mapsto(y',y'')$, where
$y'_\eps=y_{\eps0}$ and $y''_\eps=y_{\eps1}$ (the last digit is split off).
Under this identification:
\begin{itemize}
\item $L_i(y)=(L_iy',L_iy'')$ and $U_i(y)=(U_iy',U_iy'')$ for $i\le m$;
\item $L_{m+1}(y)=y'$ and $U_{m+1}(y)=y''$;
\item $y_{\vec0}=y'_{\vec0}$ and $y_{\vec1}=y''_{\vec1}$.
\end{itemize}
For a measure $\lambda$ on $X^{[m]}$ and $\eps\in[m]$, the \emph{vertex marginal} at $\eps$ is
the image of $\lambda$ under $x\mapsto x_\eps$.

\subsubsection*{Spectral measures and Wiener's lemma}
Let $\widehat\Gamma$ be the (compact metrizable) dual group of $\Gamma$; its elements are
homomorphisms $\psi\colon\Gamma\to S^1=\{z\in\mathbb C:|z|=1\}$.
For an invariant measure $\nu$, write $\mathcal U_\gamma f=f\circ T^\gamma$ for the Koopman
representation on $L^2(\nu)$.
For $f\in L^2(\nu)$, the spectral measure $\sigma_f$ is the finite positive measure on
$\widehat\Gamma$ with $\langle\mathcal U_\gamma f,f\rangle=\int\psi(\gamma)\,d\sigma_f(\psi)$
(Bochner--Herglotz); note $\sigma_f(\widehat\Gamma)=\|f\|_2^2$.
An \emph{eigenfunction with eigenvalue $\psi\in\widehat\Gamma$} is a nonzero $f$ with
$\mathcal U_\gamma f=\psi(\gamma)f$ for all $\gamma$. Write $H_\psi$ for the corresponding
eigenspace (including $0$).

\begin{lemma}\label{K:lem:wiener}
Let $\nu$ be invariant and let $\Phi$, $\Theta$ be \Fol{} sequences.
\begin{enumerate}[label=(\alph*)]
\item\label{K:wa} For $\psi\in\widehat\Gamma$, $\mathop{\mathbb E}_{\gamma\in\Phi_N}\psi(\gamma)\to1$ if
$\psi=1$, and $\to0$ otherwise.
\item\label{K:wb} For every unitary representation $V$ of $\Gamma$ on a Hilbert space and every
vector $v$, $\mathop{\mathbb E}_{\gamma\in\Phi_N}V_\gamma v$ converges to the orthogonal
projection of $v$ onto the $V$-invariant vectors.
\item\label{K:wc} For $f\in L^2(\nu)$ and $\psi\in\widehat\Gamma$, $\sigma_f(\{\psi\})=\|P_\psi f\|^2$,
where $P_\psi$ is the orthogonal projection onto $H_\psi$. In particular, $\sigma_f$ has no
atoms if and only if $f$ is orthogonal to every eigenfunction.
\item\label{K:wd} If $\sigma_f$ has no atoms, then for every $g\in L^2(\nu)$,
$\mathop{\mathbb E}_{\beta\in\Theta_M}|\langle\mathcal U_\beta f,g\rangle|\to0$ as $M\to\infty$.
\end{enumerate}
\end{lemma}

\begin{proof}
\ref{K:wa}: If $\psi(g)\ne1$, then
$|1-\psi(g)|\cdot|\mathop{\mathbb E}_{\Phi_N}\psi|\le|\Phi_N\triangle(\Phi_N+g)|/|\Phi_N|\to0$.

\ref{K:wb}: The space splits orthogonally into the invariant vectors and the closed span of
$\{u-V_gu\}$. On a vector $u-V_gu$ the averages are bounded by
$2\|u\|\,|\Phi_N\triangle(\Phi_N+g)|/|\Phi_N|$.

\ref{K:wc}: Apply \ref{K:wb} to $V_\gamma=\overline{\psi(\gamma)}\mathcal U_\gamma$, whose invariant
vectors form $H_\psi$. Then
\[
\|P_\psi f\|^2=\lim_N\mathop{\mathbb E}_{\gamma\in\Phi_N}\overline{\psi(\gamma)}\langle\mathcal U_\gamma f,f\rangle
=\lim_N\int\mathop{\mathbb E}_{\gamma\in\Phi_N}(\psi'\overline\psi)(\gamma)\,d\sigma_f(\psi')
=\sigma_f(\{\psi\})
\]
by \ref{K:wa} and dominated convergence.

\ref{K:wd}: Let $Q$ be the orthogonal projection onto the cyclic subspace $Z(f)$, the closed
span of $\{\mathcal U_\gamma f\}$.
The map $\mathcal U_\gamma f\mapsto(\psi\mapsto\psi(\gamma))$ extends to a unitary isomorphism
$Z(f)\to L^2(\sigma_f)$. Let $\phi$ be the image of $Qg$.
Then $\langle\mathcal U_\beta f,g\rangle=\langle\mathcal U_\beta f,Qg\rangle=\int\psi(\beta)\,d\tau(\psi)$, where
$\tau=\overline\phi\,\sigma_f$ is a complex measure without atoms.
By Fubini and \ref{K:wa},
\[
\mathop{\mathbb E}_{\beta\in\Theta_M}\Big|\int\psi(\beta)\,d\tau\Big|^2
=\iint\mathop{\mathbb E}_{\beta\in\Theta_M}(\psi\overline{\psi'})(\beta)\,d\tau(\psi)\,d\overline\tau(\psi')
\longrightarrow\sum_{\psi}|\tau(\{\psi\})|^2=0 .
\]
The Cauchy--Schwarz inequality finishes the proof.
\end{proof}

\subsubsection*{Kronecker maps}
Let $(Y,T)$ be a topological $\Gamma$-system and $\lambda$ an ergodic measure on $Y$.
The \emph{Kronecker subspace} $H_c(\lambda)\subset L^2(\lambda)$ is the closed linear span of
all eigenfunctions.
A \emph{rotation} is a pair $(K,\alpha)$, where $K$ is a compact metrizable abelian group
(written additively, with Haar measure $m_K$) and $\alpha\colon\Gamma\to K$ is a homomorphism
with dense image; $\Gamma$ acts on $K$ by $z\mapsto z+\alpha(\gamma)$.

\begin{definition}\label{K:def:kron}
A \emph{Kronecker map} for $(Y,\lambda)$ is a Borel map $\kappa\colon Y\to K$ to a rotation
$(K,\alpha)$ with the following properties:
\begin{itemize}
\item $\kappa_*\lambda=m_K$;
\item for every $\gamma$, $\kappa(T^\gamma y)=\kappa(y)+\alpha(\gamma)$ for $\lambda$-a.e.\ $y$;
\item $L^2(Y,\sigma(\kappa),\lambda)=H_c(\lambda)$.
\end{itemize}
We write $\mathbb E_\lambda(f\mid\kappa)$ for the orthogonal projection onto $H_c(\lambda)$,
that is, the conditional expectation onto $\sigma(\kappa)$.
\end{definition}

\begin{lemma}\label{K:lem:kron}
Let $\lambda$ be ergodic.
\begin{enumerate}[label=(\alph*)]
\item\label{K:ka} A Kronecker map for $(Y,\lambda)$ exists.
\item\label{K:kb} Let $\kappa\colon Y\to K$ be any Kronecker map, with rotation $\alpha$. Every
eigenfunction is a.e.\ equal to $c\cdot(\varphi\circ\kappa)$ for a constant $c$ and a
character $\varphi\in\widehat K$. Here $\varphi$ is the unique character with
$\varphi\circ\alpha$ equal to the eigenvalue.
\item\label{K:kc} For $f\in L^2(\lambda)$: $\mathbb E_\lambda(f\mid\kappa)=0$ if and only if
$\sigma_f$ has no atoms.
\end{enumerate}
\end{lemma}

\begin{proof}
\ref{K:ka}: This is the discrete spectrum theorem (Halmos--von Neumann; for amenable groups
see \cite[Proposition~2.12]{charamaras2024finding}). For completeness we sketch the
abelian case.
\begin{itemize}
\item The modulus of an eigenfunction is invariant, hence constant. Products of
eigenfunctions are eigenfunctions. So the set $\mathcal E\subset\widehat\Gamma$ of
eigenvalues is a subgroup.
\item Eigenfunctions with distinct eigenvalues are orthogonal and $L^2(\lambda)$ is
separable, so $\mathcal E$ is countable.
\item Let $K$ be the dual of the discrete group $\mathcal E$, which is compact and metrizable,
and let $\alpha(\gamma)=(\psi\mapsto\psi(\gamma))$. The image $\alpha(\Gamma)$ is dense, since
its annihilator in $\mathcal E$ is trivial.
\item Choose unimodular eigenfunctions $e_\psi$ for $\psi\in\mathcal E$, and set
$\Theta(y)=(e_\psi(y))_{\psi}\in(S^1)^{\mathcal E}$. Then
$\Theta(T^\gamma y)=\alpha(\gamma)\cdot\Theta(y)$ a.e., where $K$ is viewed as a closed subgroup
of $(S^1)^{\mathcal E}$.
\item Hence $\Theta_*\lambda$ is invariant under $\alpha(\Gamma)$, and so under $K$.
\item The image of $\Theta$ in $(S^1)^{\mathcal E}/K$ is invariant, so by ergodicity it is
a.e.\ a constant coset $Ku_0$.
\item Put $\kappa=\Theta\cdot u_0^{-1}$, redefined on a null set to take values in $K$.
Then $\kappa_*\lambda=m_K$ and $\kappa$ is a.e.\ equivariant.
\item By ergodicity each eigenspace is one-dimensional, spanned by $e_\psi$. The characters
of $K$ are the evaluations at elements of $\mathcal E$, so the functions $\chi\circ\kappa$,
$\chi\in\widehat K$, are constant multiples of the $e_\psi$. Hence $L^2(\sigma(\kappa))$, the
closed span of these functions, equals $H_c(\lambda)$.
\end{itemize}

\ref{K:kb}: Let $h$ be an eigenfunction with eigenvalue $\psi$. Then $h\in H_c(\lambda)$, so
$h=\tilde h\circ\kappa$ for some $\tilde h\in L^2(m_K)$, and
$\tilde h(z+\alpha(\gamma))=\psi(\gamma)\tilde h(z)$ for $m_K$-a.e.\ $z$.
Expanding $\tilde h=\sum_\varphi a_\varphi\varphi$ in characters gives
$a_\varphi(\varphi(\alpha(\gamma))-\psi(\gamma))=0$ for all $\varphi$ and $\gamma$.
Since $\alpha(\Gamma)$ is dense, at most one $\varphi$ satisfies $\varphi\circ\alpha=\psi$.

\ref{K:kc}: $H_c(\lambda)$ is the closed span of the eigenspaces $H_\psi$, so this follows
from \cref{K:lem:wiener}\ref{K:wc}.
\end{proof}

\subsubsection*{Disintegrations}
Let $\kappa\colon Y\to K$ be a Kronecker map for $(Y,\lambda)$.
There is a Borel family $z\mapsto\lambda_z$ of probability measures on $Y$ with the following
properties \cite[Theorem~2.2 and (2.1)]{charamaras2024finding}:
\begin{itemize}
\item $\lambda=\int_K\lambda_z\,dm_K(z)$;
\item $\lambda_z(\kappa^{-1}(z))=1$ for $m_K$-a.e.\ $z$;
\item for bounded Borel $f$ on $Y$,
$\mathbb E_\lambda(f\mid\kappa)=\tilde f\circ\kappa$ $\lambda$-a.e., where
\[
\tilde f(z):=\int f\,d\lambda_z .
\]
\end{itemize}
The family is unique up to $m_K$-null sets. It is a priori defined only for $m_K$-a.e.\
$z$; we extend it to all $z\in K$ by setting $\lambda_z:=\lambda$ on the exceptional null
set, which affects none of the statements below.
If $0\le f\le 1$, then $0\le\tilde f\le1$ a.e.
Since $\sigma(\kappa)$ is invariant modulo null sets, $\mathbb E_\lambda(\cdot\mid\kappa)$
commutes with $\mathcal U_\gamma$. Hence $\widetilde{f\circ T^\gamma}(z)=\tilde f(z+\alpha(\gamma))$
for $m_K$-a.e.\ $z$.

\subsection{The correspondence principle and Erdős cubes}\label{K:ssec:corr}

The following is Furstenberg's correspondence principle~\cite{furstenberg1977ergodic}, in the
form given by Charamaras and Mountakis.

\begin{theorem}[{\cite[Theorem~2.15]{charamaras2024finding}}]\label{K:thm:corr}
Let $A\subset\Gamma$, and let $\Phi$ be a \Fol{} sequence such that
$\dens_\Phi(A)=\lim_N|A\cap\Phi_N|/|\Phi_N|$ exists.
Then there exist an ergodic system $(X,\mu,T)$ on a compact metric space $X$, a clopen set
$E\subset X$, a \Fol{} sequence $\Psi$ and a point $a\in\mathrm{gen}(\mu,\Psi)$ such that
$\mu(E)\ge\dens_\Phi(A)$ and $A=\{\gamma\in\Gamma:T^\gamma a\in E\}$.
\end{theorem}

In~\cite{charamaras2024finding} this is stated for countable amenable groups. In the proof,
$X=\{0,1\}^\Gamma$ with the shift action and $a=1_A$.
We do \emph{not} need the strengthened form in which measurable eigenfunctions are
topological (property~(4) of \cite[Theorem~2.5]{kra2026short}). The re-basing lemma
(\cref{K:lem:rebase}) takes its place.

\begin{definition}[Erdős cubes; cf.\ {\cite[Definition~4.1]{kra2022infinite}}]\label{K:def:cube}
Let $(X,T)$ be a topological $\Gamma$-system and $k\ge1$.
A point $y\in X^{[k]}$ is a \emph{$k$-dimensional Erdős cube} if, for each $1\le i\le k$,
there are pairwise distinct $\beta_1,\beta_2,\ldots\in\Gamma$ with
$T^{\beta_n}L_i(y)\to U_i(y)$ in $X^{[k-1]}$.
\end{definition}

The next lemma is \cite[Lemma~4.5 and Corollary~4.6]{kra2022infinite}, transferred to
$\Gamma$. We formulate the acceptability condition so that it is meaningful while some of
the sets are still empty.

\begin{lemma}\label{K:lem:cube-sumset}
Let $y\in X^{[k]}$ be an Erdős cube, and let $E\subset X$ be open with $y_{\vec1}\in E$.
Then there are infinite sets $B_1,\dots,B_k\subset\Gamma$ such that
$T^{b_1+\dots+b_k}y_{\vec0}\in E$ for all $b_i\in B_i$.
\end{lemma}

\begin{proof}
For $\eps\in[k]$ let $Z(\eps)=\{j:\eps_j=0\}$.
Call a tuple $(B_1,\dots,B_k)$ of finite sets \emph{acceptable} if, for every $\eps\in[k]$
and every choice of $b_j\in B_j$ for $j\in Z(\eps)$,
\begin{equation}\label{K:eq:accept}
T^{\sum_{j\in Z(\eps)}b_j}\,y_\eps\in E .
\end{equation}
For $\eps=\vec1$ the sum is empty, and the condition reads $y_{\vec1}\in E$.
Hence $(\emptyset,\dots,\emptyset)$ is acceptable.

Let $(B_1,\dots,B_k)$ be acceptable, and fix $i$.
For $\eps$ with $\eps_i=0$, let
\[
H_\eps=\bigcap T^{-\sum_{j\in Z(\eps)\setminus\{i\}}b_j}E ,
\]
the intersection over all choices $b_j\in B_j$, $j\in Z(\eps)\setminus\{i\}$ (an empty
intersection is $X$).
This is a finite intersection of open sets, so $H_\eps$ is open.
Since $Z(\eps+e_i)=Z(\eps)\setminus\{i\}$, acceptability at $\eps+e_i$ says that
$y_{\eps+e_i}\in H_\eps$.

Because $y$ is an Erdős cube, there are pairwise distinct $\beta_n$ with
$T^{\beta_n}y_\eps\to y_{\eps+e_i}$ for all $\eps$ with $\eps_i=0$.
Hence for all large $n$ we have $T^{\beta_n}y_\eps\in H_\eps$ for every such $\eps$.
Choose such a $b^*=\beta_n$ with $b^*\notin B_1\cup\dots\cup B_k$, and put $B'_i=B_i\cup\{b^*\}$ and $B'_j=B_j$
for $j\ne i$.

We check~\eqref{K:eq:accept} for the new tuple. Let $\eps\in[k]$, and choose $b_j\in B'_j$
for $j\in Z(\eps)$. There are two cases.
\begin{itemize}
\item If $i\notin Z(\eps)$, or if $b_i\in B_i$, then~\eqref{K:eq:accept} is an old condition.
\item If $i\in Z(\eps)$ and $b_i=b^*$, then
$T^{\sum_{j\in Z(\eps)}b_j}y_\eps=T^{\sum_{j\in Z(\eps)\setminus\{i\}}b_j}(T^{b^*}y_\eps)\in E$,
because $T^{b^*}y_\eps\in H_\eps$. This step uses commutativity of $\Gamma$.
\end{itemize}

Now add elements in the order $i=1,2,\dots,k,1,2,\dots$.
This produces increasing acceptable tuples whose unions $B_1,\dots,B_k$ are infinite.
Each finite subfamily is contained in an acceptable stage, so~\eqref{K:eq:accept} holds for
the unions. Taking $\eps=\vec0$ gives the claim. (Since $b^*$ is always chosen outside
$B_1\cup\dots\cup B_k$, the sets $B_i$ are in addition pairwise disjoint.)
\end{proof}

\subsection{Relative products over the Kronecker factor}\label{K:ssec:relprod}

Throughout this subsection, $(Y,T)$ is a topological $\Gamma$-system, $\lambda$ is an ergodic
measure on $Y$, and $\kappa\colon Y\to K$ is a Kronecker map with rotation $\alpha$ and
disintegration $(\lambda_z)$.

\begin{definition}\label{K:def:J}
For $t\in K$, the \emph{relative product with parameter $t$} is the measure on $Y\times Y$
\[
J^\lambda_t:=\int_K\lambda_z\otimes\lambda_{z+t}\,dm_K(z).
\]
The \emph{graph measure} is $\rho_\lambda:=(\mathrm{id},\kappa)_*\lambda$ on $Y\times K$.
Here $\Gamma$ acts on $Y\times K$ by $T^\gamma\times(\cdot+\alpha(\gamma))$.
We write $\Sigma_\lambda:=\supp\rho_\lambda$.
\end{definition}

For $\Gamma=\Z$ these measures appear in \cite[(3.10) and Proposition~3.11]{kra2022infinite},
and for amenable groups in \cite[Definition~4.7 and Theorem~4.8]{charamaras2024finding}, where
$(y,y')\mapsto J^\lambda_{\kappa(y')-\kappa(y)}$ is shown to be a continuous ergodic
decomposition of $\lambda\times\lambda$ when $\kappa$ is continuous. \cref{K:lem:J} below is
the version we need; we include its short proof.

Translations preserve $m_K$-null sets. Hence $J^\lambda_t$ does not depend on the version
of the disintegration, and it is defined for \emph{every} $t\in K$.
The graph measure $\rho_\lambda$ is invariant and concentrated on the graph of $\kappa$.
The projection onto $Y$ is therefore a measure-theoretic isomorphism
$(Y\times K,\rho_\lambda)\to(Y,\lambda)$, and $\rho_\lambda$ is ergodic.
For Borel sets $S\subset Y$ and $O\subset K$ we have
\begin{equation}\label{K:eq:rho}
\rho_\lambda(S\times O)=\int_O\lambda_z(S)\,dm_K(z).
\end{equation}

\begin{lemma}\label{K:lem:J}
\leavevmode
\begin{enumerate}[label=(\alph*)]
\item\label{K:Ja} For bounded Borel $f,g$ on $Y$ and every $t\in K$,
$\int f\otimes g\,dJ^\lambda_t=\int_K\tilde f(z)\,\tilde g(z+t)\,dm_K(z)$.
\item\label{K:Jb} Each $J^\lambda_t$ is invariant under the diagonal action, has both marginals
equal to $\lambda$, and is concentrated on $\{(y,y'):\kappa(y')-\kappa(y)=t\}$.
The map $t\mapsto J^\lambda_t$ is weak$^*$ continuous.
\item\label{K:Jc} $J^\lambda_t$ is ergodic for $m_K$-almost every $t$.
\end{enumerate}
\end{lemma}

\begin{proof}
\ref{K:Ja} is the definition.

\ref{K:Jb}: The marginals are $\int\lambda_z\,dm_K=\lambda$ and $\int\lambda_{z+t}\,dm_K=\lambda$.
Concentration follows from $\lambda_z(\kappa^{-1}(z))=1$ for a.e.\ $z$.
For invariance, use \ref{K:Ja} and $\widetilde{f\circ T^\gamma}=\tilde f(\cdot+\alpha(\gamma))$:
\[
J^\lambda_t\big((f\circ T^\gamma)\otimes(g\circ T^\gamma)\big)
=\int\tilde f(z+\alpha(\gamma))\,\tilde g(z+t+\alpha(\gamma))\,dm_K(z)
=J^\lambda_t(f\otimes g).
\]
For continuity, note that $t\mapsto\int\tilde f(z)\tilde g(z+t)\,dm_K$ is continuous, by
continuity of translation in $L^2(m_K)$. Linear combinations of functions $f\otimes g$ with
$f,g\in C(Y)$ are uniformly dense in $C(Y\times Y)$.

\ref{K:Jc}: Let $D(y,y')=\kappa(y')-\kappa(y)$ and $\Lambda=\lambda\times\lambda$.

\emph{Step 1: $J$ disintegrates $\Lambda$ over $D$.}
By \ref{K:Ja} and Fubini, for bounded Borel $f,g$ on $Y$ and $h$ on $K$,
\begin{equation}\label{K:eq:disD}
\int f\otimes g\cdot(h\circ D)\,d\Lambda=\iint\tilde f(z)\,\tilde g(z')\,h(z'-z)\,dm_K\,dm_K
=\int_K h(t)\,J^\lambda_t(f\otimes g)\,dm_K(t).
\end{equation}
By a monotone class argument, \eqref{K:eq:disD} holds with $f\otimes g$ replaced by any bounded
Borel $F$ on $Y\times Y$, and $t\mapsto J^\lambda_t(F)$ is Borel.
Together with \ref{K:Jb}, this shows that $t\mapsto J^\lambda_t$ is the disintegration of
$\Lambda$ over $D$. In particular $\mathbb E_\Lambda(F\mid\sigma(D))(y,y')=J^\lambda_{D(y,y')}(F)$
for $\Lambda$-a.e.\ $(y,y')$ (recall $D_*\Lambda=m_K$).

\emph{Step 2: invariant functions of $\Lambda$ are $\sigma(D)$-measurable.}
Write $L^2(\lambda)=H_c\oplus H_w$, where $H_c=L^2(\sigma(\kappa))$ has orthonormal basis
$\{\varphi\circ\kappa:\varphi\in\widehat K\}$. Let $P_1$ be the projection onto the
invariant functions in $L^2(\Lambda)$.
\begin{itemize}
\item For $f\in H_w$ and $g\in L^2(\lambda)$, we have
$\langle\mathcal U_\gamma(f\otimes g),f\otimes g\rangle=\widehat{\sigma_f}(\gamma)\widehat{\sigma_g}(\gamma)$,
so $\sigma_{f\otimes g}=\sigma_f*\sigma_g$. This has no atoms because $\sigma_f$ has none
(\cref{K:lem:kron}\ref{K:kc}). Hence $P_1(f\otimes g)=0$ by \cref{K:lem:wiener}\ref{K:wc}.
The same holds for $g\otimes f$.
\item The function $\varphi\circ\kappa\otimes\varphi'\circ\kappa$ is an eigenfunction with
eigenvalue $(\varphi\varphi')\circ\alpha$. Since $\alpha(\Gamma)$ is dense, this eigenvalue
is trivial only if $\varphi'=\overline\varphi$.
\end{itemize}
Since $L^2(\Lambda)$ is the orthogonal sum of $H_c\otimes H_c$, $H_w\otimes L^2(\lambda)$ and
$H_c\otimes H_w$, and $P_1$ maps an eigenfunction with nontrivial eigenvalue to $0$,
$P_1L^2(\Lambda)$ is the closed span of the functions
$\varphi(\kappa(y))\overline{\varphi(\kappa(y'))}=\overline{\varphi}(D(y,y'))$, all of
which are $\sigma(D)$-measurable. Conversely $D$ is invariant.
So $P_1F=\mathbb E_\Lambda(F\mid\sigma(D))$, which is the function
$(y,y')\mapsto J^\lambda_{D(y,y')}(F)$ by Step~1.

\emph{Step 3: ergodicity of almost every $J^\lambda_t$.}
Fix a \Fol{} sequence $\Phi$, write
$A_NF=\mathop{\mathbb E}_{\gamma\in\Phi_N}F\circ T^\gamma$, and let $F\in C(Y\times Y)$.
The function $t\mapsto\|A_NF-J^\lambda_t(F)\|^2_{L^2(J^\lambda_t)}$ is continuous by \ref{K:Jb}.
Since $J^\lambda_t$ is concentrated on $D^{-1}(t)$,
\[
\int_K\big\|A_NF-J^\lambda_t(F)\big\|^2_{L^2(J^\lambda_t)}\,dm_K(t)
=\big\|A_NF-J^\lambda_{D(\cdot)}(F)\big\|^2_{L^2(\Lambda)}\xrightarrow[N\to\infty]{}0
\]
by \cref{K:lem:wiener}\ref{K:wb} and Step~2.
For each $t$, the integrand converges to $\|P^{(t)}_1F-J^\lambda_t(F)\|^2_{L^2(J^\lambda_t)}$,
where $P^{(t)}_1$ is the invariant projection in $L^2(J^\lambda_t)$.
By Fatou's lemma this limit vanishes for a.e.\ $t$.
Applying this to a countable dense subset of $C(Y\times Y)$ shows the following. For a.e.\
$t$, $P^{(t)}_1$ maps a dense subset of $L^2(J^\lambda_t)$, and hence all of it, to
constants. That is, $J^\lambda_t$ is ergodic.
\end{proof}

\begin{lemma}[Kronecker factors of factors]\label{K:lem:factor}
Let $\lambda$ be ergodic on $Y$ with Kronecker map $\kappa\colon Y\to K$ (rotation $\alpha$).
Let $P\colon Y\to Z$ be a Borel map into a topological $\Gamma$-system with
$P_*\lambda=\nu$ and, for each $\gamma$, $P\circ T^\gamma=T^\gamma\circ P$ $\lambda$-a.e., and let $\iota\colon Z\to K'$ be a Kronecker map for $\nu$ (rotation $\alpha'$).
Then:
\begin{enumerate}[label=(\alph*)]
\item\label{K:Fa} $\mathbb E_\lambda(f\circ P\mid\kappa)=\mathbb E_\nu(f\mid\iota)\circ P$ for $f\in L^2(\nu)$;
\item\label{K:Fb} there is a unique continuous homomorphism $\bar\chi\colon K\to K'$ with
$\bar\chi\circ\alpha=\alpha'$. It is surjective, and there is $c\in K'$ with
$\iota\circ P=\bar\chi\circ\kappa+c$ $\lambda$-a.e.;
\item\label{K:Fc} for \emph{every} $s\in K$, $(P\times P)_*J^\lambda_s=J^\nu_{\bar\chi(s)}$.
\end{enumerate}
\end{lemma}

\begin{proof}
\ref{K:Fa}: Write $f=f_c+f_w$ with $f_c\in H_c(\nu)$ and $f_w\perp H_c(\nu)$.
Eigenfunctions of $\nu$ pull back under $P$ to eigenfunctions of $\lambda$, so
$f_c\circ P\in H_c(\lambda)$.
Since $P$ is a factor map, $\sigma_{f_w\circ P}=\sigma_{f_w}$, which has no atoms. Hence
$f_w\circ P\perp H_c(\lambda)$ by \cref{K:lem:kron}\ref{K:kc}.

\ref{K:Fb}: Uniqueness holds because $\alpha(\Gamma)$ is dense.
For existence, take $\psi\in\widehat{K'}$. Then $\psi\circ\iota\circ P$ is an eigenfunction of
$\lambda$ with eigenvalue $\psi\circ\alpha'$.
By \cref{K:lem:kron}\ref{K:kb} there is a unique $\varphi_\psi\in\widehat K$ with
$\varphi_\psi\circ\alpha=\psi\circ\alpha'$. By uniqueness, $\psi\mapsto\varphi_\psi$ is a
homomorphism $\widehat{K'}\to\widehat K$.
By Pontryagin duality there is a continuous homomorphism $\bar\chi\colon K\to K'$ with
$\psi\circ\bar\chi=\varphi_\psi$ for all $\psi$. Then $\bar\chi\circ\alpha=\alpha'$.
The image $\bar\chi(K)$ is closed and contains the dense set $\alpha'(\Gamma)$, so $\bar\chi$
is onto.
Finally, $\iota\circ P-\bar\chi\circ\kappa$ is $\lambda$-a.e.\ invariant, since both terms
are shifted by $\alpha'(\gamma)$ under $T^\gamma$. By ergodicity it is a.e.\ equal to a
constant $c$.

\ref{K:Fc}: Let $f,g$ be bounded Borel on $Z$.
By \ref{K:Fa} and \ref{K:Fb},
$\mathbb E_\lambda(f\circ P\mid\kappa)=\tilde f(\iota\circ P)=\tilde f(\bar\chi\circ\kappa+c)$,
so $\widetilde{f\circ P}(z)=\tilde f(\bar\chi z+c)$ for $m_K$-a.e.\ $z$. Here
$\tilde f$ is formed with respect to $(\nu,\iota)$.
The surjective homomorphism $\bar\chi$ maps $m_K$ to $m_{K'}$. Hence, by \cref{K:lem:J}\ref{K:Ja},
\[
\int f\otimes g\,d(P\times P)_*J^\lambda_s=\int_K\tilde f(\bar\chi z+c)\,\tilde g(\bar\chi z+\bar\chi s+c)\,dm_K(z)
=\int_{K'}\tilde f(u)\,\tilde g(u+\bar\chi s)\,dm_{K'}(u),
\]
and the right-hand side is $J^\nu_{\bar\chi(s)}(f\otimes g)$.
Both sides are probability measures, so they agree.
\end{proof}

The next lemma is \cite[Lemma~3.18]{kra2022infinite} for $\Gamma=\Z$, and the Claim in the
proof of \cite[Theorem~4.10]{charamaras2024finding} for amenable groups. We include a short
proof, based on the van der Corput inequality rather than on the continuity of the ergodic
decomposition. In the argument of~\cite{kra2026short} its role is played by the choice of a
generic $y$ such that $(x,y)$ is generic for some (not necessarily ergodic) joining.

\begin{lemma}[Generic pairs]\label{K:lem:pairs}
Suppose in addition that $\pi\colon Y\to K$ is a \emph{continuous} Kronecker map with
$\pi(T^\gamma y)=\pi(y)+\alpha(\gamma)$ for \emph{all} $y$ and $\gamma$. Let $(J_t)$ be the
relative products formed with respect to $\pi$.
Let $y_0\in\mathrm{gen}(\lambda,\Phi)$.
Then there are a subsequence $\Phi'$ of $\Phi$ and a Borel set $Y_0\subset Y$ with
$\lambda(Y_0)=1$ such that
\[
(y_0,y)\in\mathrm{gen}\big(J_{\pi(y)-\pi(y_0)},\Phi'\big)\qquad\text{for every }y\in Y_0 .
\]
\end{lemma}

\begin{proof}
Let $f,g\in C(Y)$.
We show that, in $L^2(\lambda)$ (as functions of $y$),
\begin{equation}\label{K:eq:L2conv}
\mathop{\mathbb E}_{\gamma\in\Phi_N}f(T^\gamma y_0)\,g(T^\gamma y)\ \longrightarrow\
\Lambda_{f,g}(y):=J_{\pi(y)-\pi(y_0)}(f\otimes g)=\int\tilde f(z)\,\tilde g(z+\pi(y)-\pi(y_0))\,dm_K(z).
\end{equation}
Fix $\tau>0$. Choose a trigonometric polynomial $\tilde f'=\sum_jc_j\varphi_j$ on $K$
with $\|\tilde f-\tilde f'\|_{L^2(m_K)}<\tau$, and put $f'=\tilde f'\circ\pi$.
The function $f'$ is continuous because $\pi$ is.
Let $f''=f-f'=f^\perp+f^\sharp$, with $f^\perp=f-\mathbb E_\lambda(f\mid\pi)$ and
$f^\sharp=(\tilde f-\tilde f')\circ\pi$. Then $\|f^\sharp\|_2<\tau$, and $\sigma_{f^\perp}$ has no atoms
(\cref{K:lem:kron}\ref{K:kc}).

\emph{The part $f'$.}
We have $f'(T^\gamma y_0)=\sum_jc_j\varphi_j(\pi y_0)\,\varphi_j(\alpha(\gamma))$.
By \cref{K:lem:wiener}\ref{K:wb} applied to $V_\gamma=\varphi_j(\alpha(\gamma))\mathcal U_\gamma$, the
average $\mathop{\mathbb E}_{\Phi_N}\varphi_j(\alpha(\gamma))\mathcal U_\gamma g$ converges to the
projection of $g$ onto $H_{\overline{\varphi_j\circ\alpha}}$. That eigenspace is spanned by
$\overline{\varphi_j}\circ\pi$, so the projection is
$\big(\int\tilde g\,\varphi_j\,dm_K\big)\,\overline{\varphi_j}\circ\pi$.
Hence the averages for $f'$ converge in $L^2(\lambda)$ to
\[
\Lambda'_{f,g}(y)=\sum_jc_j\,\varphi_j(\pi y_0)\,\overline{\varphi_j(\pi y)}\int\tilde g\,\varphi_j\,dm_K .
\]
A change of variables shows that $\Lambda'_{f,g}(y)=\int\tilde f'(z)\tilde g(z+\pi(y)-\pi(y_0))\,dm_K(z)$.
It follows that $|\Lambda_{f,g}-\Lambda'_{f,g}|\le\tau\|g\|_\infty$ everywhere.

\emph{The part $f''$.}
Put $u(\gamma)=f''(T^\gamma y_0)$. By the spectral theorem,
\[
\Big\|\mathop{\mathbb E}_{\gamma\in\Phi_N}u(\gamma)\mathcal U_\gamma g\Big\|_2^2=\int_{\widehat\Gamma}\Big|\mathop{\mathbb E}_{\gamma\in\Phi_N}u(\gamma)\psi(\gamma)\Big|^2d\sigma_g(\psi).
\]
We bound the integrand uniformly in $\psi$ by a van der Corput argument.
Fix a finite $H\subset\Gamma$. By the \Fol{} property,
$\mathop{\mathbb E}_{\Phi_N}u\psi=\mathop{\mathbb E}_{\gamma\in\Phi_N}\mathop{\mathbb E}_{h\in H}(u\psi)(\gamma+h)+o(1)$.
By Cauchy--Schwarz,
\[
\Big|\mathop{\mathbb E}_{\Phi_N}u\psi\Big|^2\le\frac1{|H|^2}\sum_{h,h'\in H}\Big|\mathop{\mathbb E}_{\gamma\in\Phi_N}f''(T^{\gamma+h}y_0)\overline{f''(T^{\gamma+h'}y_0)}\Big|+o(1).
\]
The function $f''$ is continuous and $y_0$ is generic, so the inner averages converge to
$\langle\mathcal U_{h-h'}f'',f''\rangle$.
Moreover,
$|\langle\mathcal U_\delta f'',f''\rangle-\langle\mathcal U_\delta f^\perp,f^\perp\rangle|\le\tau(\|f^\perp\|_2+\|f''\|_2)\le\tau(2\|f\|_\infty+\tau)$.
Take $H=\Theta_M$ for a \Fol{} sequence $\Theta$. By Cauchy--Schwarz and the computation
in \cref{K:lem:wiener}\ref{K:wd},
\[
\frac1{|\Theta_M|^2}\sum_{h,h'\in\Theta_M}|\langle\mathcal U_{h-h'}f^\perp,f^\perp\rangle|\le\Big(\iint\Big|\mathop{\mathbb E}_{h\in\Theta_M}(\psi\overline{\psi'})(h)\Big|^2d\sigma_{f^\perp}\,d\sigma_{f^\perp}\Big)^{1/2}\to0 .
\]
Hence $\sup_\psi\limsup_N|\mathop{\mathbb E}_{\Phi_N}u\psi|^2\le\tau(2\|f\|_\infty+\tau)$.
The integrands are bounded and $\sigma_g$ is finite, so by reverse Fatou
\[
\limsup_N\Big\|\mathop{\mathbb E}_{\Phi_N}u(\gamma)\mathcal U_\gamma g\Big\|_2^2\le\tau(2\|f\|_\infty+\tau)\|g\|_2^2 .
\]
Combining the two parts,
\[
\limsup_N\Big\|\mathop{\mathbb E}_{\gamma\in\Phi_N}f(T^\gamma y_0)\,\mathcal U_\gamma g-\Lambda_{f,g}\Big\|_2
\le\tau\|g\|_\infty+\big(\tau(2\|f\|_\infty+\tau)\big)^{1/2}\|g\|_2 ,
\]
and letting $\tau\to0$ proves~\eqref{K:eq:L2conv}.

Now let $\mathcal D$ be a countable dense subset of $C(Y)$.
Convergence in $L^2$ implies a.e.\ convergence along a subsequence. A diagonal argument
gives a subsequence $\Phi'$ and a conull Borel set $Y_0$ such that the convergence
in~\eqref{K:eq:L2conv} holds at every $y\in Y_0$ along $\Phi'$, for all $f,g\in\mathcal D$.
Both sides of~\eqref{K:eq:L2conv} are bounded bilinear in $(f,g)$ with respect to the
uniform norm. The span of $\{f\otimes g\}$ is dense in $C(Y\times Y)$.
Hence $(y_0,y)\in\mathrm{gen}(J_{\pi(y)-\pi(y_0)},\Phi')$ for $y\in Y_0$.
\end{proof}

\subsection{Re-basing and two support lemmas}\label{K:ssec:support}

The following lemma is the rotation case of \cite[Lemma~3.5]{charamaras2024finding} (see
the remark following that lemma, which allows any topologically distal factor in place of
the Kronecker factor). Charamaras and Mountakis prove it for amenable groups, extending a
result of Host and Kra for $\Z$-actions \cite[Proposition~6.1]{host2009uniformity} (see also
\cite[Proposition~3.19]{kra2022infinite}); their proof is in
\cite[Appendix~A]{charamaras2024finding}.
For rotations the Ellis-semigroup step of that proof reduces to compactness of $K$, so we
include the short argument.

\begin{lemma}[Re-basing]\label{K:lem:rebase}
Let $(W,T)$ be a topological $\Gamma$-system and let $w\in W$ have dense orbit.
Let $\lambda$ be an ergodic measure on $W$.
Let $\kappa\colon W\to K$ be a Borel map into a compact metrizable abelian group, and
$\alpha\colon\Gamma\to K$ a homomorphism, such that for each $\gamma$,
$\kappa(T^\gamma x)=\kappa(x)+\alpha(\gamma)$ for $\lambda$-a.e.\ $x$.
Let $\hat\lambda=(\mathrm{id},\kappa)_*\lambda$ on $W\times K$, with the action
$T^\gamma\times(\cdot+\alpha(\gamma))$.
Then there are $d\in K$ and a \Fol{} sequence $\Psi$ such that $(w,d)\in\mathrm{gen}(\hat\lambda,\Psi)$.
\end{lemma}

\begin{proof}
The measure $\hat\lambda$ is invariant and ergodic, since it is isomorphic to $\lambda$ via the
projection.
By \ref{K:G4} there are a point $(x_1,z_1)$ and a \Fol{} sequence $\Phi$ with
$(x_1,z_1)\in\mathrm{gen}(\hat\lambda,\Phi)$.

For each $N$, the finitely many maps $T^\gamma$, $\gamma\in\Phi_N$, are continuous at $x_1$,
and the orbit of $w$ is dense. So there is $h_N\in\Gamma$ with
$\sup_{\gamma\in\Phi_N}d_W(T^{\gamma+h_N}w,T^\gamma x_1)<1/N$.

By compactness, $\alpha(h_{N_j})\to u_*$ in $K$ along a subsequence. Put $d=z_1-u_*$ and
$\Psi_j=\Phi_{N_j}+h_{N_j}$; this is again a \Fol{} sequence.
For $F\in C(W\times K)$,
\[
\mathop{\mathbb E}_{\gamma\in\Psi_j}F(T^\gamma w,d+\alpha(\gamma))
=\mathop{\mathbb E}_{\gamma\in\Phi_{N_j}}F\big(T^{\gamma+h_{N_j}}w,\ z_1+\alpha(\gamma)+(\alpha(h_{N_j})-u_*)\big).
\]
The arguments are uniformly within $o(1)$ of $(T^\gamma x_1,z_1+\alpha(\gamma))$: in the first
coordinate by the choice of $h_N$, and in the second by uniform continuity of addition on
$K$. Uniform continuity of $F$ and genericity of $(x_1,z_1)$ give the claim.
\end{proof}

\begin{lemma}[Support of the graph measure]\label{K:lem:stab}
Let $(Z,T)$ be a topological $\Gamma$-system, $\nu$ ergodic on $Z$, and $\iota\colon Z\to K$ a
Kronecker map with rotation $\alpha$. Let $\rho=\rho_\nu$ and $\Sigma=\supp\rho\subset Z\times K$.
Suppose $(p,\theta)\in Z\times K$ is generic for $\rho$ along some \Fol{} sequence, and
$p\in\supp\nu$. Then $(p,\theta)\in\Sigma$.
\end{lemma}

\begin{proof}
The projection of $\Sigma$ to $Z$ is compact and has full $\nu$-measure, so it contains
$\supp\nu$. Pick $\theta_0$ with $(p,\theta_0)\in\Sigma$ and put $\zeta=\theta-\theta_0$.

For $u\in K$ let $S_u(x,t)=(x,t+u)$. These maps commute with the action, so $S_\zeta\Sigma$
is closed and invariant. It contains $S_\zeta(p,\theta_0)=(p,\theta)$, and hence the orbit closure
of $(p,\theta)$, which contains $\Sigma$ by~\ref{K:G2}.
So $\Sigma\subset S_\zeta\Sigma$, that is, $S_{-\zeta}\Sigma\subset\Sigma$, and hence
$S_{-n\zeta}\Sigma\subset\Sigma$ for all $n\ge1$.

By compactness there are $1\le n_1<n_2<\cdots$ with $n_{l+1}-n_l\ge2$ and $n_l\zeta\to\zeta_*$.
Then $(n_{l+1}-n_l)\zeta\to0$, so with $q_l=n_{l+1}-n_l-1\ge1$ we get $-q_l\zeta\to\zeta$.
Since $\Sigma$ is closed, $S_\zeta\Sigma\subset\Sigma$, and so
$(p,\theta)=S_\zeta(p,\theta_0)\in\Sigma$.
\end{proof}

\begin{lemma}[Phase returns]\label{K:lem:phase}
Let $(W,T)$ and $(Z,T)$ be topological $\Gamma$-systems, $\lambda$ ergodic on $W$, and
$w\in\mathrm{gen}(\lambda,\Psi)$ for some \Fol{} sequence $\Psi$.
Let $L,U\colon W\to Z$ be continuous equivariant maps with $L_*\lambda=U_*\lambda=:\nu$.
Let $\iota\colon Z\to K$ be a Kronecker map for $\nu$ with rotation $\alpha$.
\begin{enumerate}[label=(\alph*)]
\item\label{K:Pa} There is $r\in K$ with $\iota\circ U=\iota\circ L+r$ $\lambda$-a.e.
\item\label{K:Pb} If $U(w)\in\supp\nu$, then for every open $V\ni U(w)$ and every open $O\ni0$ in $K$,
there is a \Fol{} sequence $\Omega$ along which the set
\[
\{\gamma\in\Gamma:\ T^\gamma L(w)\in V,\ \alpha(\gamma)\in r+O\}
\]
has positive lower density; in particular this set has positive upper Banach density.
\end{enumerate}
\end{lemma}

\begin{proof}
\ref{K:Pa}: Both $\iota\circ U$ and $\iota\circ L$ are shifted by $\alpha(\gamma)$ under
$T^\gamma$ (a.e.), so their difference is $\lambda$-a.e.\ invariant, hence a.e.\ constant.

\ref{K:Pb}: Let $W'$ be the orbit closure of $w$. Then $\lambda(W')=1$ by~\ref{K:G2}, and $w$
has dense orbit in $W'$.
The map $\kappa'=\iota\circ L$ satisfies the hypotheses of \cref{K:lem:rebase} on $W'$.
That lemma gives $z\in K$ and a \Fol{} sequence $\Omega$ with $(w,z)$ generic along $\Omega$
for $\hat\lambda=(\mathrm{id},\iota\circ L)_*\lambda$.
Consider the continuous equivariant maps $(x,t)\mapsto(Lx,t)$ and $(x,t)\mapsto(Ux,t+r)$
from $W'\times K$ to $Z\times K$.
They push $\hat\lambda$ forward to $(L,\iota\circ L)_*\lambda=\rho_\nu$ and to
$(U,\iota\circ L+r)_*\lambda=(U,\iota\circ U)_*\lambda=\rho_\nu$, respectively.
By~\ref{K:G3}, both $(L(w),z)$ and $(U(w),z+r)$ are generic for $\rho_\nu$ along $\Omega$.
By \cref{K:lem:stab}, $(U(w),z+r)\in\supp\rho_\nu$.
So $V\times(z+r+O)$ is an open set meeting $\supp\rho_\nu$, and by~\ref{K:G1} the set
\[
\{\gamma:(T^\gamma L(w),z+\alpha(\gamma))\in V\times(z+r+O)\}=\{\gamma:T^\gamma L(w)\in V,\ \alpha(\gamma)\in r+O\}
\]
has positive lower density along $\Omega$.
\end{proof}

\begin{lemma}[Pointwise support]\label{K:lem:PS}
Let $\nu$, $\iota\colon Z\to K$, $\rho=\rho_\nu$, $\Sigma=\supp\rho$ and $J_t=J^\nu_t$ be as above.
If $(p,\theta)\in\Sigma$, then for $\rho$-almost every $(q,\varphi)\in Z\times K$,
\[
(p,q)\in\supp J_{\varphi-\theta}.
\]
More precisely, the set of $(q,\varphi)$ with $(p,q)\notin\supp J_{\varphi-\theta}$ is
contained in a $\rho$-null Borel set.
\end{lemma}

\begin{proof}
Fix open sets $O_1\ni p$ and $O_2\subset Z$.
Define Borel functions $f(\xi)=\nu_\xi(O_1)$ and $g(\xi)=\nu_\xi(O_2)$ on $K$, with values in
$[0,1]$, and let $F(t)=J_t(O_1\times O_2)=\int f(\xi)g(\xi+t)\,dm_K(\xi)$.
The function $F$ is continuous.

By~\eqref{K:eq:rho}, $\int_Nf\,dm_K=\rho(O_1\times N)>0$ for every open $N\ni\theta$, because
$O_1\times N$ is a neighbourhood of $(p,\theta)\in\Sigma$.

Let $C=\{\varphi\in K:F(\varphi-\theta)=0\}$, a closed set. We claim $\int_Cg\,dm_K=0$.
Suppose not.
Put $R(k)=\int_Cg(\varphi)\,g(k+\varphi-\theta)\,dm_K(\varphi)$.
Then $R\ge0$, and $R$ is continuous since
$|R(k)-R(k')|\le\|g(\cdot+k)-g(\cdot+k')\|_{L^1(m_K)}$.
Moreover $R(\theta)=\int_Cg^2\,dm_K>0$. So $R>R(\theta)/2$ on some open $N\ni\theta$, and
\[
\int f\,R\,dm_K\ \ge\ \tfrac12R(\theta)\int_Nf\,dm_K>0 .
\]
On the other hand, by Fubini,
\[
\begin{aligned}
\int f(k)R(k)\,dm_K(k)&=\int_Cg(\varphi)\int f(k)\,g(k+\varphi-\theta)\,dm_K(k)\,dm_K(\varphi)\\
&=\int_Cg(\varphi)F(\varphi-\theta)\,dm_K(\varphi)=0,
\end{aligned}
\]
a contradiction.

By~\eqref{K:eq:rho}, $\int_Cg\,dm_K=\rho\big(O_2\times C\big)$. So
\[
\mathcal N_{O_1,O_2}:=\{(q,\varphi):\ q\in O_2,\ J_{\varphi-\theta}(O_1\times O_2)=0\}=O_2\times C
\]
is a $\rho$-null Borel set.
Let $\mathcal B$ be a countable base of $Z$.
If $(p,q)\notin\supp J_{\varphi-\theta}$, then some $O_1\times O_2$ with $O_1,O_2\in\mathcal B$,
$p\in O_1$ and $q\in O_2$ satisfies $J_{\varphi-\theta}(O_1\times O_2)=0$, so
$(q,\varphi)\in\mathcal N_{O_1,O_2}$. The union of these countably many $\rho$-null sets is
$\rho$-null.
\end{proof}

\begin{remark}\label{K:rem:PS}
For a fixed point $p$, the proof uses only that $\theta\in\supp(f_{O_1}m_K)$ for every open
$O_1\ni p$, where $f_{O_1}(\xi)=\nu_\xi(O_1)$. By~\eqref{K:eq:rho} this is equivalent to
$(p,\theta)\in\Sigma$.
\end{remark}

\subsection{Cube data and the cube criterion}\label{K:ssec:criterion}

From now on, $(X,\mu,T)$ is an ergodic system on a compact metric space, $a\in X$, and
$E\subset X$ is open with $\mu(E)>0$.

\begin{definition}\label{K:def:datum}
Let $m\ge0$. An \emph{$m$-dimensional cube datum} is a triple $(w,\lambda,\Psi)$ with the
following properties:
\begin{enumerate}[label=(D\arabic*)]
\item\label{K:D1} $\lambda$ is an ergodic measure on $X^{[m]}$ (for the diagonal action),
$\Psi$ is a \Fol{} sequence, $w\in\mathrm{gen}(\lambda,\Psi)$, and $w_{\vec0}=a$;
\item\label{K:D2} every vertex marginal of $\lambda$ equals $\mu$;
\item\label{K:D3} (face symmetry) $(L_i)_*\lambda=(U_i)_*\lambda=:\nu_i$ for $1\le i\le m$;
\item\label{K:D4} (upper faces in the support) $U_i(w)\in\supp\nu_i$ for $1\le i\le m$.
\end{enumerate}
\end{definition}

If $a\in\mathrm{gen}(\mu,\Psi)$, then $(a,\mu,\Psi)$ is a $0$-dimensional cube datum, since
\ref{K:D3} and \ref{K:D4} are empty.
Each face measure $\nu_i$ is ergodic, being a factor of $\lambda$. In what follows,
$\iota_i\colon X^{[m-1]}\to K_i$ denotes a Kronecker map for $\nu_i$, with rotation $\alpha_i$.

The next proposition is our version of Claim~1 in the proof of
\cite[Theorem~3.1]{kra2026short}, together with the iteration that follows it.
The structure of the argument comes from there:
\begin{itemize}
\item the splitting into a Kronecker part and a weak mixing part;
\item the positivity remark \cite[Remark~2.2]{kra2026short};
\item the averaging of the weak mixing part along \emph{every} \Fol{} sequence.
\end{itemize}
What is new is the treatment of the Kronecker part. We use continuity of translation on
the Kronecker group of the face measure and the phase returns of \cref{K:lem:phase}. No
continuity of eigenfunctions is needed.

\begin{proposition}[Cube criterion]\label{K:prop:criterion}
If an $m$-dimensional cube datum $(w,\lambda,\Psi)$ exists, then there is $x\in X^{[m]}$ such
that $y=(w,x)\in X^{[m+1]}$ is an $(m+1)$-dimensional Erdős cube with $y_{\vec0}=a$ and
$y_{\vec1}=x_{\vec1}\in E$.
\end{proposition}

\begin{proof}
Fix a metric on $X$ and use the maximum metric on products.
An \emph{open box} is a set $V=\prod_{\eps\in[m]}V_\eps\subset X^{[m]}$ with each $V_\eps$
open.
For $1\le i\le m$ and $\beta\in\Gamma$ put $U_i(V)=\prod_{\eta\in[m-1]}V_{\eta^{(i,1)}}$ and
\[
W_i(V,\beta)=\{x\in V:\ T^\beta L_i(x)\in U_i(V)\}
=\{x\in V:\ T^\beta x_{\eta^{(i,0)}}\in V_{\eta^{(i,1)}}\ \forall\eta\in[m-1]\}.
\]
This set is open, and $V\subset U_i^{-1}(U_i(V))$.

\smallskip\noindent\emph{Claim A (direction $m+1$).} If $V$ is open with $\lambda(V)>0$, there are
infinitely many $\beta$ with $T^\beta w\in V$.

This is~\ref{K:G1} together with~\ref{K:D1}.

\smallskip\noindent\emph{Claim B (directions $i\le m$).} If $V$ is an open box with $\lambda(V)>0$,
$1\le i\le m$ and $\delta>0$, then there are infinitely many $\beta$ such that
\[
d\big(T^\beta L_i(w),U_i(w)\big)<\delta\qquad\text{and}\qquad\lambda\big(W_i(V,\beta)\big)>0 .
\]

\emph{Proof of Claim B.}
Let $F=1_{U_i(V)}$ on $X^{[m-1]}$. Since $L_i$ is equivariant,
\begin{equation}\label{K:eq:Wi}
\lambda(W_i(V,\beta))=\int 1_V(x)\,F(L_iT^\beta x)\,d\lambda(x)=\langle F\circ L_i\circ T^\beta,1_V\rangle_\lambda .
\end{equation}
In $L^2(\nu_i)$, write $F=F_c+F_w$, where $F_c=\mathbb E_{\nu_i}(F\mid\iota_i)=G\circ\iota_i$
with $G\colon K_i\to[0,1]$ Borel, and $F_w\perp H_c(\nu_i)$.
Let $r_i$ be the constant of \cref{K:lem:phase}\ref{K:Pa} for $L=L_i$ and $U=U_i$, so that
$\iota_i\circ U_i=\iota_i\circ L_i+r_i$ $\lambda$-a.e.

\emph{Positivity.}
Let $\Delta=\{F_c>0\}$, a $\sigma(\iota_i)$-measurable set. Then
$\int_{\Delta^c}F\,d\nu_i=\int_{\Delta^c}F_c\,d\nu_i=0$, so $\nu_i(U_i(V)\setminus\Delta)=0$
\cite[Remark~2.2]{kra2026short}.
Since $(U_i)_*\lambda=\nu_i$ and $V\subset U_i^{-1}(U_i(V))$, we get $F_c(U_ix)>0$ for
$\lambda$-a.e.\ $x\in V$. Hence
\[
\eta_V:=\int_VF_c(U_ix)\,d\lambda(x)>0 .
\]

\emph{The Kronecker part.}
For every $\beta$ and $\lambda$-a.e.\ $x$ we have
$F_c(T^\beta L_ix)=G(\iota_i(L_ix)+\alpha_i(\beta))$ and
$F_c(U_ix)=G(\iota_i(L_ix)+r_i)$. Here we use that $L_i$ and $U_i$ push $\lambda$ to $\nu_i$,
that $\iota_i$ is a.e.\ equivariant, and the choice of $r_i$.
Since $(\iota_i\circ L_i)_*\lambda=m_{K_i}$,
\[
\big|\langle F_c\circ L_i\circ T^\beta,1_V\rangle_\lambda-\eta_V\big|
\le\int_{K_i}\big|G(z+\alpha_i(\beta))-G(z+r_i)\big|\,dm_{K_i}(z)
=\big\|G(\cdot+\alpha_i(\beta)-r_i)-G\big\|_{L^1(m_{K_i})}.
\]
By continuity of translation in $L^1(m_{K_i})$, there is an open $O\ni0$ in $K_i$ such that
the right-hand side is $<\eta_V/4$ whenever $\alpha_i(\beta)\in r_i+O$.

\emph{The weak mixing part.}
Let $h=F_w\circ L_i\in L^2(\lambda)$, so that $F_w\circ L_i\circ T^\beta=\mathcal U_\beta h$.
Since $L_i$ is a factor map,
$\langle\mathcal U_\beta h,h\rangle_\lambda=\langle\mathcal U_\beta F_w,F_w\rangle_{\nu_i}$, so
$\sigma_h=\sigma_{F_w}$, which has no atoms by \cref{K:lem:kron}\ref{K:kc}.
By \cref{K:lem:wiener}\ref{K:wd}, for every \Fol{} sequence $\Theta$,
\[
\mathop{\mathbb E}_{\beta\in\Theta_M}|\langle\mathcal U_\beta h,1_V\rangle_\lambda|\to0 .
\]

\emph{Assembly.}
Let $\mathcal N_\delta$ be the open $\delta$-ball around $U_i(w)$, and let
\[
R_\delta=\{\beta:\ T^\beta L_i(w)\in\mathcal N_\delta,\ \alpha_i(\beta)\in r_i+O\}.
\]
By~\ref{K:D4} and \cref{K:lem:phase}\ref{K:Pb}, applied with $L=L_i$, $U=U_i$ and
$\nu=\nu_i$, there are a \Fol{} sequence $\Omega$ and $c_0>0$ with
$|R_\delta\cap\Omega_M|\ge c_0|\Omega_M|$ for all large $M$.
Let $R'_\delta=\{\beta\in R_\delta:|\langle\mathcal U_\beta h,1_V\rangle|\ge\eta_V/4\}$. By the weak
mixing part,
$|R'_\delta\cap\Omega_M|\le(4/\eta_V)\sum_{\beta\in\Omega_M}|\langle\mathcal U_\beta h,1_V\rangle|=o(|\Omega_M|)$.
Hence $|(R_\delta\setminus R'_\delta)\cap\Omega_M|\ge(c_0/2)|\Omega_M|$ for all large $M$, and
$R_\delta\setminus R'_\delta$ is infinite by \cref{K:lem:folner-grow}.
For $\beta\in R_\delta\setminus R'_\delta$, by~\eqref{K:eq:Wi},
\[
\lambda(W_i(V,\beta))=\langle F_c\circ L_i\circ T^\beta,1_V\rangle+\langle\mathcal U_\beta h,1_V\rangle
\ge\eta_V-\tfrac{\eta_V}4-\tfrac{\eta_V}4>0 .
\]
This proves Claim~B.

\smallskip\noindent\emph{Nested boxes.}
Every open subset of $X^{[m]}$ is a countable union of open boxes whose closures lie in it
and whose coordinates have arbitrarily small diameter. So every open set of positive
$\lambda$-measure contains such a box of positive measure.
Let $V_0=\{x\in X^{[m]}:x_{\vec1}\in E\}$. This is an open box, and
$\lambda(V_0)=\mu(E)>0$ by~\ref{K:D2}.
Let $i_n\in\{1,\dots,m+1\}$ with $i_n\equiv n\pmod{m+1}$.

Suppose an open box $V_n$ with $\lambda(V_n)>0$ has been chosen. Choose
$\beta_n\notin\{\beta_0,\dots,\beta_{n-1}\}$ as follows.
\begin{itemize}
\item If $i_n=m+1$, let $\beta_n$ satisfy $T^{\beta_n}w\in V_n$ (Claim~A), and put $W^*=V_n$.
\item If $i_n=i\le m$, let $\beta_n$ satisfy $d(T^{\beta_n}L_i(w),U_i(w))<1/(n+1)$ and
$\lambda(W_i(V_n,\beta_n))>0$ (Claim~B), and put $W^*=W_i(V_n,\beta_n)$.
\end{itemize}
Then choose an open box $V_{n+1}$ with $\overline{V_{n+1}}\subset W^*$, $\lambda(V_{n+1})>0$, and all
coordinate diameters $<1/(n+1)$.

The compact sets $\overline{V_n}$ are nested and their diameters tend to $0$. So
$\bigcap_n\overline{V_n}=\{x\}$, and $x\in V_n$ for all $n$; in particular $x_{\vec1}\in E$.
Let $y=(w,x)$.
\begin{itemize}
\item If $i_n=m+1$, then $T^{\beta_n}w$ and $x$ both lie in $V_n$, so
$d(T^{\beta_n}L_{m+1}(y),U_{m+1}(y))=d(T^{\beta_n}w,x)\le\operatorname{diam}V_n$.
\item If $i_n=i\le m$, then $x\in W_i(V_n,\beta_n)$. So for every $\eta\in[m-1]$,
both $T^{\beta_n}x_{\eta^{(i,0)}}$ and $x_{\eta^{(i,1)}}$ lie in $V_{n,\eta^{(i,1)}}$, and
\[
d\big(T^{\beta_n}L_i(x),U_i(x)\big)\le\operatorname{diam}V_n .
\]
Together with $d(T^{\beta_n}L_i(w),U_i(w))<1/(n+1)$, and $L_i(y)=(L_iw,L_ix)$,
$U_i(y)=(U_iw,U_ix)$, this gives $T^{\beta_n}L_i(y)\to U_i(y)$ along these $n$.
\end{itemize}
The $\beta_n$ are pairwise distinct, so $y$ is an Erdős cube. Moreover $y_{\vec0}=w_{\vec0}=a$
and $y_{\vec1}=x_{\vec1}\in E$.
\end{proof}

\subsection{Building the cube level by level}\label{K:ssec:induction}

\begin{proposition}[Induction step]\label{K:prop:step}
If an $m$-dimensional cube datum exists, then an $(m+1)$-dimensional cube datum exists.
\end{proposition}

\begin{proof}
Let $(w,\lambda,\Psi)$ be an $m$-dimensional cube datum, with face measures $\nu_i$ and
Kronecker maps $\iota_i\colon X^{[m-1]}\to K_i$ (rotations $\alpha_i$), $1\le i\le m$.
Let $\kappa\colon X^{[m]}\to K$ be a Kronecker map for $\lambda$, with rotation $\alpha$.
For $i\le m$, apply \cref{K:lem:factor} to $P=U_i$ and to $P=L_i$. Both push $\lambda$ to
$\nu_i$, so they give the same homomorphism $\bar\chi_i\colon K\to K_i$, the unique one with
$\bar\chi_i\circ\alpha=\alpha_i$, together with constants $c_i^U,c_i^L\in K_i$ such that
\begin{equation}\label{K:eq:affine}
\iota_i\circ U_i=\bar\chi_i\circ\kappa+c_i^U,\qquad\iota_i\circ L_i=\bar\chi_i\circ\kappa+c_i^L\qquad\lambda\text{-a.e.}
\end{equation}
Moreover, for every $s\in K$,
\begin{equation}\label{K:eq:facesym}
(U_i\times U_i)_*J^\lambda_s=(L_i\times L_i)_*J^\lambda_s=J^{\nu_i}_{\bar\chi_i(s)} .
\end{equation}

\emph{Step 1: re-basing.}
This is the graph construction of \cite[Proposition~3.20]{kra2022infinite} and
\cite[Proposition~3.7]{charamaras2024finding}, applied to the current measure $\lambda$.
Let $W$ be the orbit closure of $w$ in $X^{[m]}$, so that $\lambda(W)=1$ by~\ref{K:G2}.
By \cref{K:lem:rebase} there are $d\in K$ and a \Fol{} sequence $\Psi'$ with
$(w,d)\in\mathrm{gen}(\hat\lambda,\Psi')$, where $\hat\lambda=(\mathrm{id},\kappa)_*\lambda$.
Let $\widehat W\subset W\times K$ be the orbit closure of $(w,d)$; then $\hat\lambda(\widehat W)=1$.
The projection $\mathrm{pr}\colon\widehat W\to X^{[m]}$ is a measure-theoretic isomorphism
from $\hat\lambda$ to $\lambda$, with inverse $x\mapsto(x,\kappa(x))$.
So $\pi(x,t)=t$ is a \emph{continuous} Kronecker map for $\hat\lambda$, it is equivariant
everywhere, and $\kappa\circ\mathrm{pr}=\pi$ $\hat\lambda$-a.e.
By \cref{K:lem:factor}\ref{K:Fc} (with $\bar\chi=\mathrm{id}$ and $c=0$),
\begin{equation}\label{K:eq:prJ}
(\mathrm{pr}\times\mathrm{pr})_*J^{\hat\lambda}_s=J^\lambda_s\qquad\text{for every }s\in K .
\end{equation}

\emph{Step 2: the new top face.}
Apply \cref{K:lem:pairs} to $(\widehat W,\hat\lambda,\pi)$ and $y_0=(w,d)$.
This gives a subsequence $\Psi''$ of $\Psi'$ and a $\hat\lambda$-conull set $\widehat W_0$
such that, for $\hat x\in\widehat W_0$,
\[
((w,d),\hat x)\in\mathrm{gen}\big(J^{\hat\lambda}_{\pi(\hat x)-d},\Psi''\big).
\]
For $x\in X^{[m]}$ put $\hat x=(x,\kappa(x))$ and $s(x)=\kappa(x)-d$.
Since $\hat\lambda=(\mathrm{id},\kappa)_*\lambda$, we have $\hat x\in\widehat W_0$ for
$\lambda$-a.e.\ $x$.
For such $x$, projecting with~\ref{K:G3} and~\eqref{K:eq:prJ} gives
\begin{equation}\label{K:eq:newgen}
w':=(w,x)\in\mathrm{gen}\big(\lambda':=J^\lambda_{s(x)},\ \Psi''\big)\qquad\text{in }X^{[m]}\times X^{[m]}=X^{[m+1]} .
\end{equation}
Note that $s(x)$ is $m_K$-distributed when $x\sim\lambda$.
By \cref{K:lem:J}\ref{K:Jc}, $\lambda'=J^\lambda_{s(x)}$ is ergodic for $\lambda$-a.e.\ $x$.

\emph{Step 3: \ref{K:D1}--\ref{K:D3} for $(w',\lambda',\Psi'')$.}
\begin{itemize}
\item \ref{K:D1}: $\lambda'$ is invariant and (for a.e.\ $x$) ergodic, $w'$ is generic by
\eqref{K:eq:newgen}, and $w'_{\vec0}=w_{\vec0}=a$.
\item \ref{K:D2}: by \cref{K:lem:J}\ref{K:Jb} both marginals of $\lambda'$ on $X^{[m]}$ equal
$\lambda$. So the vertex marginals of $\lambda'$ are those of $\lambda$, all equal to $\mu$.
\item \ref{K:D3}, direction $m+1$: $(L_{m+1})_*\lambda'$ and $(U_{m+1})_*\lambda'$ are the two
marginals of $\lambda'$, both equal to $\lambda$.
\item \ref{K:D3}, direction $i\le m$: $L_i(y',y'')=(L_iy',L_iy'')$, and similarly for $U_i$.
By~\eqref{K:eq:facesym}, $(L_i)_*\lambda'=(U_i)_*\lambda'=J^{\nu_i}_{t_i}$ with
$t_i:=\bar\chi_i(s(x))$.
\end{itemize}

\emph{Step 4: \ref{K:D4} for $(w',\lambda',\Psi'')$.}

\emph{Direction $m+1$.} We need $U_{m+1}(w')=x\in\supp\lambda$, which holds for
$\lambda$-a.e.\ $x$.

\emph{Direction $i\le m$.}
Let $p=U_i(w)$ and $\theta=\bar\chi_i(d)+c_i^U$.
The map $(z,u)\mapsto(U_iz,\bar\chi_i(u)+c_i^U)$ from $X^{[m]}\times K$ to $X^{[m-1]}\times K_i$
is continuous. It is equivariant because $\bar\chi_i\circ\alpha=\alpha_i$.
By~\eqref{K:eq:affine} it pushes $\hat\lambda$ to
$(U_i,\iota_i\circ U_i)_*\lambda=(\mathrm{id},\iota_i)_*\nu_i=\rho_{\nu_i}$.
By~\ref{K:G3}, $(p,\theta)$ is generic for $\rho_{\nu_i}$ along $\Psi'$.
Since $p\in\supp\nu_i$ by~\ref{K:D4}, \cref{K:lem:stab} gives $(p,\theta)\in\supp\rho_{\nu_i}$.
Now let
\[
(q,\varphi):=\big(U_ix,\ \iota_i(U_ix)\big).
\]
When $x\sim\lambda$, the pair $(q,\varphi)$ is distributed according to $\rho_{\nu_i}$.
By~\eqref{K:eq:affine}, for $\lambda$-a.e.\ $x$,
\[
\varphi-\theta=\bar\chi_i(\kappa(x))-\bar\chi_i(d)=\bar\chi_i(s(x))=t_i .
\]
By \cref{K:lem:PS}, applied to $(p,\theta)\in\supp\rho_{\nu_i}$, for $\lambda$-a.e.\ $x$ we have
\[
U_i(w')=(U_iw,U_ix)=(p,q)\in\supp J^{\nu_i}_{\varphi-\theta}=\supp J^{\nu_i}_{t_i}=\supp(U_i)_*\lambda' .
\]

\emph{Conclusion.}
The requirements on $x$ in Steps 2--4 are finitely many: $\hat x\in\widehat W_0$, ergodicity of
$J^\lambda_{s(x)}$, $x\in\supp\lambda$, and the $m$ conditions from \cref{K:lem:PS}.
Each holds for $\lambda$-almost every $x$.
The data they depend on, namely $d$, $\Psi'$, $\Psi''$, $\widehat W_0$ and the points
$(p,\theta)$, were all fixed before $x$ was chosen.
So some $x$ satisfies all of them, and $(w',\lambda',\Psi'')$ is an $(m+1)$-dimensional
cube datum.
\end{proof}


\subsection{Proof of \texorpdfstring{\cref{K:thm:main}}{the main theorem}}\label{K:ssec:proof}

Let $A\subset\Gamma$ with $\ubd(A)>0$, and let $k\ge1$.
Choose a \Fol{} sequence along which $|A\cap\Phi_N|/|\Phi_N|$ has a positive limsup, and pass
to a subsequence along which the limit $\dens_\Phi(A)>0$ exists.
\cref{K:thm:corr} gives an ergodic $(X,\mu,T)$, a clopen $E$ with $\mu(E)>0$, a \Fol{}
sequence $\Psi$, and $a\in\mathrm{gen}(\mu,\Psi)$ with $A=\{\gamma:T^\gamma a\in E\}$.
\begin{enumerate}
\item $(a,\mu,\Psi)$ is a $0$-dimensional cube datum.
\item By \cref{K:prop:step}, applied $k-1$ times, there is a $(k-1)$-dimensional cube datum.
\item By \cref{K:prop:criterion}, there is a $k$-dimensional Erdős cube $y$ with $y_{\vec0}=a$
and $y_{\vec1}\in E$.
\item By \cref{K:lem:cube-sumset}, there are infinite $B_1,\dots,B_k$ with
$T^{b_1+\dots+b_k}a\in E$, that is, $b_1+\dots+b_k\in A$, for all $b_i\in B_i$.
\end{enumerate}
\qed

\formalization{\Cref{K:thm:main} is formalized in the fourteen files
\nolinkurl{Kmrr25/ThmK*.lean}; the main statement is \leanname{thmK} in \leanfile{ThmK}. All
lemmas and propositions of this section are proved there, with four classical results stated as
explicit hypotheses (the structure \leanname{ClassicalInputs} in \leanfile{ThmKDefs}): the
correspondence principle (\cref{K:thm:corr}); the existence of generic points~\ref{K:G4}; the
discrete spectrum theorem (\cref{K:lem:kron}\ref{K:ka}); and, for vectors orthogonal to all
eigenfunctions, the decay of averaged correlations given by the spectral theorem and Wiener's
lemma (\cref{K:lem:wiener}\ref{K:wc},\ref{K:wd}). Disintegrations are taken from Mathlib. Some
arguments are organised differently in the formal proof: for instance, the Kronecker group is
realised as a closed subgroup of $(\R/\Z)^\N$, which avoids Pontryagin duality, and
\cref{K:lem:pairs} splits $g$ into its Kronecker and weak mixing parts instead of using the
spectral theorem for $g$.}

\subsection{Comments on the proof}\label{K:ssec:comments}

\subsubsection*{Structure of the induction}
The only information passed from one level to the next is the cube datum
\ref{K:D1}--\ref{K:D4}. The decoration $d$ and the Kronecker maps are used within a single
step and are rebuilt for the new measure. The face measures $J^{\nu_i}_{t_i}$ of the new
measure may have much larger Kronecker factors than the $\nu_i$, but nothing about them is
needed: at the next level they enter only through \ref{K:D3} and \ref{K:D4}.
The measures $\lambda_m$ need not be invariant under face transformations or permutations
of digits, since the cube criterion treats each direction separately. Each direction
produces its own sequence of return times, so the method gives $k$ unrelated sets $B_i$.
It says nothing about sumsets $B\oplus B+t$ built from a single set, which can fail when
$\{g+g:g\in G\}$ has infinite index (see the remark after Conjecture~5.15
in~\cite{kra2023problems}).

\subsubsection*{Ingredients}
Besides elementary facts about \Fol{} sequences, the proof uses:
\begin{itemize}
\item Furstenberg's correspondence principle;
\item the existence of generic points for ergodic measures (via Lindenstrauss' pointwise
theorem);
\item the mean ergodic theorem, Wiener's lemma and the van der Corput inequality;
\item the discrete spectrum theorem for the Kronecker factor and Pontryagin duality;
\item disintegration of measures;
\item re-basing for rotations.
\end{itemize}
It uses no characteristic factors beyond the Kronecker factor, no Host--Kra or
Conze--Lesigne structure theory, no nilsystems, and no equidistribution theorem for
nilsystems. Nor does
it use finite generation, or any assumption on the subgroups $n\Gamma$.
Commutativity of $\Gamma$ is used throughout: in the spectral theory on $\widehat\Gamma$, in
describing the Kronecker factor as a rotation of a compact abelian group, and in the fact
that translates $\Phi_N+h_N$ of a \Fol{} sequence are again \Fol{}. It enters the
combinatorics in \cref{K:lem:cube-sumset} and in the diagonal action on faces.
Infiniteness of $\Gamma$ is used to obtain infinitely many distinct return times
(\cref{K:lem:folner-grow}).

\subsubsection*{Relation to earlier work}
The relative products, their ergodicity, the generic-pairs lemma and the graph construction
used for re-basing are known in this context
\cite{kra2022infinite,charamaras2024finding}; see the references given with
\cref{K:def:J}, \cref{K:lem:pairs,K:lem:rebase} and Step~1 of \cref{K:prop:step}.
For $k=2$ our argument differs from that of~\cite{kra2026short} in three respects:
re-basing replaces topological eigenfunctions, phase returns
(\cref{K:lem:phase,K:lem:stab}) replace the equicontinuity estimate~(4.2)
of~\cite{kra2026short}, and an ergodic relative product $J^\mu_s$ replaces an unspecified
joining.
The ingredients that are new relative to these arguments are the treatment of the Kronecker part of the face measures through
phase returns (\cref{K:prop:criterion}), the support lemmas
(\cref{K:lem:stab,K:lem:phase,K:lem:PS}), and the level-by-level induction
(\cref{K:prop:step}).
